\documentclass[UTF8,zihao=-4,fontset=windows]{ctexart}
\usepackage[a4paper,margin=2.25cm,headheight=15pt]{geometry}
\usepackage{graphicx,booktabs,tabularx,array,enumitem,xcolor,fancyhdr,titlesec,caption,hyperref,float,amsmath,listings,placeins,needspace}
\usepackage{colortbl}
\definecolor{pale}{HTML}{F3F3F3}
\definecolor{codebg}{HTML}{F6F6F6}
\hypersetup{colorlinks=true,linkcolor=black,urlcolor=black,citecolor=black,pdfborder={0 0 0},
  pdftitle={自动驾驶技术第三次实验报告},pdfauthor={杜承敖}}
\setlength{\parindent}{2em}
\setlength{\parskip}{0.28em}
\setlength{\emergencystretch}{2em}
\renewcommand{\arraystretch}{1.18}
\setcounter{tocdepth}{2}
\captionsetup{font=small,labelfont={bf},skip=6pt}
\setlist[itemize]{leftmargin=2em,itemsep=0.2em,topsep=0.3em}
\lstset{language=Python,basicstyle=\ttfamily\footnotesize\color{black},frame=single,
  backgroundcolor=\color{codebg},rulecolor=\color{black},columns=fullflexible,
  breaklines=true,breakatwhitespace=true,keepspaces=true,showstringspaces=false,
  keywordstyle=\color{black},commentstyle=\color{black},stringstyle=\color{black},
  numberstyle=\color{black},xleftmargin=0.6em,xrightmargin=0.6em,
  aboveskip=0.45em,belowskip=0.35em,captionpos=b}
\renewcommand{\lstlistingname}{代码清单}
\titleformat{\section}{\Large\bfseries\heiti\color{black}}{\thesection}{0.7em}{}
\titleformat{\subsection}{\large\bfseries\heiti\color{black}}{\thesubsection}{0.7em}{}
\titlespacing*{\section}{0pt}{1.0em}{0.45em}
\titlespacing*{\subsection}{0pt}{0.7em}{0.3em}
\pagestyle{fancy}
\fancyhf{}
\fancyhead[L]{\small 自动驾驶技术第三次实验}
\fancyhead[R]{\small nuScenes mini 与 PointPillars 三维检测}
\fancyfoot[C]{\small\color{black}\thepage}
\renewcommand{\headrulewidth}{0.4pt}
\widowpenalty=10000
\clubpenalty=10000
\begin{document}
\color{black}
\begin{titlepage}
  \thispagestyle{empty}
  \begin{center}
    \vspace*{1.8cm}
    {\zihao{1}\bfseries\heiti 自动驾驶技术第三次实验报告\par}
    \vspace{1.0cm}
    {\zihao{3}\bfseries nuScenes mini\par}
    \vspace{0.25cm}
    {\zihao{3}\bfseries LiDAR BEV 三维目标检测与 PointPillars\par}
    \vfill
    {\zihao{4}\renewcommand{\arraystretch}{1.75}
    \begin{tabular}{>{\bfseries}r l}
      姓名： & \underline{\makebox[9.2cm][c]{杜承敖}} \\
      学号： & \underline{\makebox[9.2cm][c]{2024302112019}} \\
      专业： & \underline{\makebox[9.2cm][c]{计算机科学与技术（雷军班）}} \\
      课程： & \underline{\makebox[9.2cm][c]{自动驾驶技术}} \\
    \end{tabular}}
    \vfill
    {\zihao{3}\bfseries 武汉大学\par}
    \vspace{0.45cm}
    {\zihao{4}2026 年 10 月 13 日\par}
    \vspace*{0.8cm}
  \end{center}
\end{titlepage}
\pagenumbering{arabic}
\tableofcontents
\clearpage
\phantomsection
\section*{摘要}
\addcontentsline{toc}{section}{摘要}

本报告记录基于 PointPillars 的 LiDAR BEV 三维目标检测实践。实验使用 nuScenes mini 的 404 个关键帧，按官方场景划分为 323 帧训练数据和 81 帧验证数据。首先比较点云与图像的表达方式、优势和互补关系；随后结合 HardVFE、BEV 映射、二维骨干网络及锚框检测头源码理解处理流程。在 RTX 4060 Laptop GPU 上，从公开权重开始微调检测头，完成 1 轮训练后恢复至累计 5 轮，共 1615 个批次。整轮平均损失由 2.873274 降至 2.456259。使用同一验证集，NDS 从 0.026860 提升至 0.118891，mAP 从 0.000061 提升至 0.062302。报告包含真实训练 Log、Loss 曲线、参数更新审计和三帧 BEV 预测图，并说明教学范围与 mini 数据的局限。


\noindent\textbf{关键词：}nuScenes mini；LiDAR；PointPillars；柱特征；三维目标检测

\clearpage

\Needspace{8\baselineskip}
\section{实验目的与数据}

按课程指导书，先评估公开 PointPillars 权重，再进行检测头微调：训练 1 轮后从检查点恢复至累计 5 轮，最后在同一 mini 验证集比较指标，并显示 LiDAR 点云及预测框。


复用实验二的 nuScenes mini 原始 JSON 元数据，并从实验一的本地点云压缩包提取对应的全部 404 个关键帧。重新生成 MMDetection3D 1.4 格式的 LiDAR 索引，未使用 PETR 模型、相机索引或预测结果，也未下载额外数据集。


\begin{table}[H]
\centering\small
\caption{运行环境与数据设置}
\label{tab:1}
\begin{tabularx}{\textwidth}{>{\raggedright\arraybackslash}p{0.30723\dimexpr\textwidth-4\tabcolsep\relax}>{\raggedright\arraybackslash}X}
\toprule
\rowcolor{pale}
\textbf{项目} & \textbf{实际设置} \\
\midrule
{GPU} & {NVIDIA GeForce RTX 4060 Laptop GPU} \\
{Python} & {3.9.25} \\
{PyTorch /\allowbreak{} CUDA} & {2.0.1+cu118 /\allowbreak{} 11.8} \\
{MMCV /\allowbreak{} MMEngine} & {2.1.0 /\allowbreak{} 0.10.7} \\
{MMDetection /\allowbreak{} MMDetection3D} & {3.2.0 /\allowbreak{} 1.4.0} \\
{数据划分} & {8 个训练场景 /\allowbreak{} 323 帧；2 个验证场景 /\allowbreak{} 81 帧；无交叉} \\
{输入} & {单 LiDAR 关键帧；原始五维存储，模型读取 x/\allowbreak{}y/\allowbreak{}z/\allowbreak{}intensity} \\
{感知区域 /\allowbreak{} BEV} & {x、y 为 [-25,25] 米；z 为 [-5,3] 米；200 × 200 网格} \\
\bottomrule
\end{tabularx}
\end{table}
404 个点云文件均存在且数据有限；标注采用 LiDAR 坐标、长度/\allowbreak{}宽度/\allowbreak{}高度和 yaw，并保存 LiDAR 坐标速度。所有三维框中心回转至全局坐标后的最大误差为 4.547e-13 米。


未读取历史 sweeps 或相机图像。数据按官方 mini\_train/\allowbreak{}mini\_val 场景划分；验证集不完整覆盖全部十个检测类别。



\FloatBarrier

\Needspace{8\baselineskip}
\section{LiDAR 点云与图像：区别、优势及互补}

点云和图像描述同一交通场景，但信息组织方式不同。图像在像素网格上记录外观，LiDAR 根据激光回波测得空间采样点。点云中的 intensity 是回波强度，不能直接当作相机的 RGB 颜色。


\begin{table}[H]
\centering\small
\caption{LiDAR 点云与图像数据的比较}
\label{tab:2}
\begin{tabularx}{\textwidth}{>{\raggedright\arraybackslash}p{0.16265\dimexpr\textwidth-6\tabcolsep\relax}>{\raggedright\arraybackslash}p{0.41867\dimexpr\textwidth-6\tabcolsep\relax}>{\raggedright\arraybackslash}X}
\toprule
\rowcolor{pale}
\textbf{比较维度} & \textbf{LiDAR 点云} & \textbf{相机图像} \\
\midrule
{数据表达} & {稀疏、无固定点序的三维点集合；本次输入为 x/\allowbreak{}y/\allowbreak{}z/\allowbreak{}intensity。} & {规则的二维像素网格，通常含 RGB 三通道，细节密集。} \\
{几何与距离} & {直接提供距离和三维位置，便于确定物体尺寸、朝向和障碍物空间。} & {像素主要表示投影；单幅普通 RGB 图像不能直接给出唯一的绝对深度。} \\
{外观与语义} & {能反映三维轮廓、反射强度，但颜色、文字和细纹理信息有限。} & {颜色、纹理、交通灯状态、标志文字和目标外观信息丰富。} \\
{采样与遮挡} & {远处及小目标往往点数较少，遮挡后可见表面采样更少。} & {远处目标仍可能有可辨外观，但像素变少、遮挡也会影响识别。} \\
{光照与天气} & {主动测距，对环境可见光变化相对不敏感；雨雾、反射特性和遮挡仍可影响回波。} & {成本相对低、图像分辨率高；弱光、眩光、雨雾和曝光变化会降低可用信息。} \\
\bottomrule
\end{tabularx}
\end{table}
各自优势：LiDAR 更适合提供有物理尺度的空间几何，图像更适合提供密集外观和细粒度语义。例如，点云有助于判断前车距离与占据空间，相机有助于辨认车辆外观和交通灯颜色。任何一种传感器都不能仅凭自身保证覆盖所有困难场景。


互补方法：先统一时间戳和坐标系，用外参把点云变换到相机坐标，再用内参投影到图像。在相机前方且可见的点处采样图像特征或语义类别概率，与点的几何特征拼接，随后进行三维检测；也可分别检测后在统一坐标中关联目标。


\Needspace{6\baselineskip}
\begin{lstlisting}[caption={投影关系的简化表达：还需处理畸变、遮挡和不同采集时刻的车辆运动。}]
p_camera = T_camera_from_lidar @ [x, y, z, 1]
u = fx * X_camera / Z_camera + cx
v = fy * Y_camera / Z_camera + cy
# Require Z_camera > 0 and a valid visible image position.
\end{lstlisting}

{\small PointPainting [2] 是这种互补思路的例子：将图像分割类别分数附加到对应点云，再送入点云检测器。本次课程模型仅使用 LiDAR，未实现相机融合；此处说明的是可扩展的融合原理。\par}

{\small\noindent [2] Vora 等，PointPainting，CVPR 2020。 \href{https://arxiv.org/abs/1911.10150}{原文链接}\par}


\FloatBarrier

\Needspace{8\baselineskip}
\section{PointPillars 原理与代码理解}

PointPillars 的关键思路是把无规则点云组织成竖直柱，学习每柱的特征，再映射为规则 BEV 伪图像，使后续网络可以使用高效的二维卷积 [1]。本实验采用课程提供的 HardVFE 编码实现；以下解释以实际配置和源码为准。


以下源码路径相对于 third\_party/\allowbreak{}mmdetection3d-1.4.0/\allowbreak{}mmdet3d/\allowbreak{}models/\allowbreak{}；课程配置与冻结策略分别见 course\_common.py 和 course\_modules.py。


\Needspace{8\baselineskip}
\subsection{从点云输入到柱特征}

course\_common.py 的 single\_frame\_pipeline() 使用 LoadPointsFromFile(load\_dim=5, use\_dim=4)：原始文件按五维点记录解析，送入模型的是前四维。训练阶段同步对点云和真实框旋转、缩放、翻转，并过滤范围和类别；验证阶段没有随机增强，也没有加载历史 sweeps。


柱划分：voxel\_size=[0.25,0.25,8]，感知范围为 50×50×8 米，因此 x/\allowbreak{}y 各 200 格，z 只有一格。虽然框架称此操作为 voxelization，本配置得到的是贯穿高度范围的柱。每柱最多 64 点；训练最多保留 6000 个柱，验证最多 8000 个，空缺点通过补零组织成张量。


\begin{table}[H]
\centering\small
\caption{柱特征的组织形式}
\label{tab:3}
\begin{tabularx}{\textwidth}{>{\raggedright\arraybackslash}p{0.24699\dimexpr\textwidth-4\tabcolsep\relax}>{\raggedright\arraybackslash}X}
\toprule
\rowcolor{pale}
\textbf{符号 /\allowbreak{} 设置} & \textbf{含义} \\
\midrule
{K} & {该帧实际保留的非空柱数量，随点云变化。} \\
{[K,64,4]} & {柱中点的原始输入组织形式；另有 num\_points 保存真实点数。} \\
{10 维点特征} & {4 维输入 + 到柱内点均值的 3 维偏移 + 到柱中心的 3 维偏移。} \\
{[K,64]} & {最终每个柱的 64 维学习特征。} \\
\bottomrule
\end{tabularx}
\end{table}
\Needspace{6\baselineskip}
\begin{lstlisting}[caption={源码摘录：HardVFE.forward()，voxel\_encoders/\allowbreak{}voxel\_encoder.py 第 409-414 行。}]
points_mean = (
    features[:, :, :3].sum(dim=1, keepdim=True) /
    num_points.type_as(features).view(-1, 1, 1))
# TODO: maybe also do cluster for reflectivity
f_cluster = features[:, :, :3] - points_mean
features_ls.append(f_cluster)
\end{lstlisting}

上述均值只除以真实点数，f\_cluster 描述柱内局部形状；f\_center 描述点相对于固定网格中心的位置。拼接后的补零位置由 get\_paddings\_indicator() 生成的 mask 清零，防止把补齐位置当作观测点。高度 z 仍保留在特征中，使用 BEV 并不表示完全丢掉高度。


HardVFE 中两层 VFELayer 使用共享线性变换、BatchNorm 和 ReLU。第一层把点特征与该柱的最大池化特征拼接，第二层再池化得到一个 64 维柱向量；这将不同点数的柱变为等长表示。本配置的特征装饰是 10 维，不能套用其他实现的固定维数。


{\small 固定验证帧 0 的实际前向记录：范围过滤后输入 32636 个四维点，形成 3550 个非空柱。第一层输入为 [3550, 64, 10]，输出为 [3550, 64, 128]；第二层输出为 [3550, 64]。\par}

{\small\noindent [1] Lang 等，PointPillars，CVPR 2019。 \href{https://openaccess.thecvf.com/content_CVPR_2019/html/Lang_PointPillars_Fast_Encoders_for_Object_Detection_From_Point_Clouds_CVPR_2019_paper.html}{原文链接}\par}


\FloatBarrier

\Needspace{8\baselineskip}
\subsection{BEV 映射、二维网络与检测头}

\Needspace{6\baselineskip}
\begin{lstlisting}[caption={源码摘录：MVXTwoStageDetector.extract\_pts\_feat()，detectors/\allowbreak{}mvx\_two\_stage.py 第 207-219 行。}]
if not self.with_pts_bbox:
    return None
voxel_features = self.pts_voxel_encoder(voxel_dict['voxels'],
                                        voxel_dict['num_points'],
                                        voxel_dict['coors'], img_feats,
                                        batch_input_metas)
batch_size = voxel_dict['coors'][-1, 0] + 1
x = self.pts_middle_encoder(voxel_features, voxel_dict['coors'],
                            batch_size)
x = self.pts_backbone(x)
if self.with_pts_neck:
    x = self.pts_neck(x)
return x
\end{lstlisting}

这段调用链依次完成柱编码、BEV 映射、骨干特征提取和多尺度融合。PointPillarsScatter 为每帧创建全零画布，柱坐标按 [batch,z,y,x] 保存；使用 y*nx+x 计算平面位置，将对应柱特征写入画布。空柱维持零值，恢复为 [B,64,200,200]，这里的通道是学习特征，并非 RGB。


\Needspace{6\baselineskip}
\begin{lstlisting}[caption={源码摘录：PointPillarsScatter.forward\_batch()，middle\_encoders/\allowbreak{}pillar\_scatter.py 第 85-93 行。}]
batch_mask = coors[:, 0] == batch_itt
this_coors = coors[batch_mask, :]
indices = this_coors[:, 2] * self.nx + this_coors[:, 3]
indices = indices.type(torch.long)
voxels = voxel_features[batch_mask, :]
voxels = voxels.t()

# Now scatter the blob back to the canvas.
canvas[:, indices] = voxels
\end{lstlisting}

SECOND 骨干用二维卷积逐级降低空间分辨率、扩大感受野；SECONDFPN 对三个尺度分别处理并对齐到同一网格，再拼接为 384 通道。Anchor3DHead 用三个 1×1 卷积分支输出类别、三维框编码和方向分类，解码后通过旋转 NMS 去除重复框。


\begin{table}[H]
\centering\small
\caption{固定验证帧 0 的中间特征尺寸}
\label{tab:4}
\begin{tabularx}{\textwidth}{>{\raggedright\arraybackslash}p{0.30723\dimexpr\textwidth-4\tabcolsep\relax}>{\raggedright\arraybackslash}X}
\toprule
\rowcolor{pale}
\textbf{实际单帧前向记录} & \textbf{输出尺寸} \\
\midrule
{HardVFE} & {[3550, 64]} \\
{PointPillarsScatter} & {[1, 64, 200, 200]} \\
{SECOND 三尺度} & {[[1, 64, 100, 100], [1, 128, 50, 50], [1, 256, 25, 25]]} \\
{SECONDFPN} & {[[1, 384, 100, 100]]} \\
{分类分支} & {[1, 140, 100, 100]} \\
{框回归分支} & {[1, 126, 100, 100]} \\
{方向分支} & {[1, 28, 100, 100]} \\
\bottomrule
\end{tabularx}
\end{table}
{\small 上表是在第 5 轮权重上对固定验证帧 0 额外执行一次 eval 前向得到的，未更新模型。每个网格位置有 14 个锚框；分类维度为锚框数×10，回归维度为锚框数×9（含位置、尺寸、yaw、vx/\allowbreak{}vy 编码），方向维度为锚框数×2。完整记录见 model\_shape\_trace.json。\par}

{\small 代码中的 loss\_cls 为 FocalLoss，loss\_bbox 为 SmoothL1Loss，loss\_dir 为方向交叉熵（权重 0.2）；训练 Log 的 loss 是已加权损失之和。course\_modules.py 按参数名仅打开 pts\_bbox\_head.* 的梯度，所以三个卷积分支的 weight/\allowbreak{}bias 共 6 个张量参加参数更新。冻结参数不代表 BatchNorm 的运行统计缓冲区也固定。\par}


\FloatBarrier

\Needspace{8\baselineskip}
\section{训练设置、断点恢复与 Loss 曲线}

\subsection{训练设置与损失曲线}

从已校验 SHA256 的公开 f19d00a3 权重初始化。仅检测头参数参与优化；Pillar 编码器、骨干及 neck 参数冻结。单 GPU、batch=1、AMP、AdamW，学习率 0.0001，权重衰减 0.01，梯度裁剪 max\_norm=35，随机种子 20261012。


\begin{table}[H]
\centering\small
\caption{首轮训练和断点续训审计}
\label{tab:5}
\begin{tabularx}{\textwidth}{>{\raggedright\arraybackslash}p{0.18072\dimexpr\textwidth-8\tabcolsep\relax}>{\raggedright\arraybackslash}p{0.27711\dimexpr\textwidth-8\tabcolsep\relax}>{\raggedright\arraybackslash}p{0.21084\dimexpr\textwidth-8\tabcolsep\relax}>{\raggedright\arraybackslash}X}
\toprule
\rowcolor{pale}
\textbf{运行} & \textbf{完成轮次 /\allowbreak{} 批次数} & \textbf{训练审计耗时} & \textbf{参数更新 /\allowbreak{} 梯度张量} \\
\midrule
{首次训练} & {第 1 轮 /\allowbreak{} 323} & {20.50 秒} & {6/\allowbreak{}6；6} \\
{断点续训} & {第 2-5 轮 /\allowbreak{} 1292} & {56.39 秒} & {6/\allowbreak{}6；6} \\
\bottomrule
\end{tabularx}
\end{table}
两段审计均为 PARAMETERS\_UPDATED，存在真实梯度和参数更新。两段运行的检查点分别确认完成第 1 轮和累计第 5 轮；检查全部模型浮点权重及 1615 个训练损失点均为有限值。续训恢复优化器及训练进度，--epochs 5 表示累计五轮。


\begin{figure}[H]
\centering
\includegraphics[width=0.98\textwidth]{\detokenize{images/loss_all_5epochs.png}}
\caption{五轮累计损失。虚线为第 1 轮后断点恢复位置；橙色点为每轮全部 323 个批次损失均值。}
\label{fig:1}
\end{figure}
\begin{table}[H]
\centering\small
\caption{各轮全部批次的平均训练损失}
\label{tab:6}
\begin{tabularx}{\textwidth}{>{\raggedright\arraybackslash}p{0.19277\dimexpr\textwidth-12\tabcolsep\relax}>{\raggedright\arraybackslash}p{0.16145\dimexpr\textwidth-12\tabcolsep\relax}>{\raggedright\arraybackslash}p{0.16145\dimexpr\textwidth-12\tabcolsep\relax}>{\raggedright\arraybackslash}p{0.16145\dimexpr\textwidth-12\tabcolsep\relax}>{\raggedright\arraybackslash}p{0.16145\dimexpr\textwidth-12\tabcolsep\relax}>{\raggedright\arraybackslash}X}
\toprule
\rowcolor{pale}
\textbf{轮次} & \textbf{1} & \textbf{2} & \textbf{3} & \textbf{4} & \textbf{5} \\
\midrule
{平均损失} & {2.87327} & {2.68935} & {2.58483} & {2.51116} & {2.45626} \\
\bottomrule
\end{tabularx}
\end{table}
{\small 损失下降反映训练优化，不等同于验证精度提升。第 1 轮及续训损失曲线、训练日志、梯度审计和检查点均单独保留。训练审计耗时包含该训练运行内部的验证等开销；逐步命令耗时另见 execution.json。\par}

{\small AMP 审计记录 3 次非有限缩放梯度事件；该计数属于动态损失缩放诊断，不能表述为所有原始梯度均有限。本次最终参数和训练损失检查通过。\par}


\FloatBarrier

\Needspace{8\baselineskip}
\subsection{训练过程 Log：实际运行摘录}

以下内容直接来自本次保存的控制台日志，页面仅自动折行，未重新构造数值。完整日志位于 outputs/\allowbreak{}setup/\allowbreak{}run1\_console.log 和 run5\_console.log，训练目录内另有框架日志及结构化记录。


\Needspace{6\baselineskip}
\begin{lstlisting}[caption={首轮训练：第 10 和第 320 个批次。}]
10/13 20:00:58 - mmengine - INFO - Epoch(train) [1][ 10/323]  lr: 1.0000e-04  eta: 0:00:38  time: 0.1220  data_time: 0.0133  memory: 295  grad_norm: 2.7219  loss: 3.2840  loss_cls: 1.1094  loss_bbox: 1.9927  loss_dir: 0.1818
10/13 20:01:11 - mmengine - INFO - Epoch(train) [1][320/323]  lr: 1.0000e-04  eta: 0:00:00  time: 0.0419  data_time: 0.0095  memory: 313  grad_norm: 2.4906  loss: 2.5351  loss_cls: 0.9419  loss_bbox: 1.4118  loss_dir: 0.1814
\end{lstlisting}

\Needspace{6\baselineskip}
\begin{lstlisting}[caption={断点恢复：恢复到 epoch=1、iter=323，继续第 2 轮并完成第 5 轮。}]
10/13 20:01:33 - mmengine - INFO - resumed epoch: 1, iter: 323
10/13 20:01:33 - mmengine - INFO - Epoch(train) [2][ 10/323]  lr: 1.0000e-04  eta: 0:01:32  time: 0.0476  data_time: 0.0082  memory: 279  grad_norm: 2.4806  loss: 2.5871  loss_cls: 0.9739  loss_bbox: 1.4223  loss_dir: 0.1910
10/13 20:02:21 - mmengine - INFO - Epoch(train) [5][320/323]  lr: 1.0000e-04  eta: 0:00:00  time: 0.0352  data_time: 0.0038  memory: 294  grad_norm: 2.7527  loss: 2.5562  loss_cls: 0.8707  loss_bbox: 1.5123  loss_dir: 0.1732
\end{lstlisting}

\Needspace{6\baselineskip}
\begin{lstlisting}[caption={两段训练的参数更新审计原文。}]
10/13 20:01:17 - mmengine - INFO - AUDIT: PARAMETERS_UPDATED; changed=6; max_change=2.187e-02
10/13 20:02:29 - mmengine - INFO - AUDIT: PARAMETERS_UPDATED; changed=6; max_change=8.097e-02
\end{lstlisting}

Log 读法：Epoch(train) [N][i/\allowbreak{}323] 表示第 N 轮、第 i 个批次；lr 是学习率，time/\allowbreak{}data\_time 是运行记录中的迭代/\allowbreak{}数据时间，memory 是日志记录的显存值；loss\_cls、loss\_bbox、loss\_dir 分别对应分类、框回归和方向损失。日志中的损失汇总由框架日志处理器输出，单条批次记录不能替代整轮均值。


断点恢复不是重新初始化模型：train.py 将 resume=True 与 load\_from 指向 epoch\_1.pth，Runner 恢复权重、优化器和累计迭代状态。日志明确出现 resumed epoch: 1, iter: 323，而后出现第 2 轮记录；最终检查点元数据为 epoch=5，证明累计训练五轮。


全部批次记录的第 1 轮平均 loss 为 2.87327，第 5 轮为 2.45626，下降约 14.51\%。各批次受点数、目标分布及数据增强影响而波动，因此结合完整曲线和轮均值观察训练。3 次缩放梯度非有限事件已保留在审计中，最终权重及损失有限性检查通过。


\Needspace{6\baselineskip}
\begin{lstlisting}[caption={实际训练、断点续训与最终验证命令。}]
python train.py --epochs 1 --work-dir outputs/run1_compact --amp
python train.py --resume-from outputs/run1_compact/epoch_1.pth --epochs 5 --work-dir outputs/run5_compact --amp
python test.py --checkpoint outputs/run5_compact/epoch_5.pth --output-dir outputs/finetuned_compact --vis-count 3 --score-thr 0.10
\end{lstlisting}


\FloatBarrier

\Needspace{8\baselineskip}
\section{验证结果与阈值解释}

公开权重和第 5 轮权重均对同一清单的全部 81 个 mini\_val 帧完成推理与 nuScenes 检测评估。根据预测中的样本索引和点云路径回查，81 个 token 及顺序均一致；mini 无独立带标注测试集，本报告的测试为验证集评估。


\begin{table}[H]
\centering\small
\caption{微调前后 NDS 与 mAP}
\label{tab:7}
\begin{tabularx}{\textwidth}{>{\raggedright\arraybackslash}p{0.19277\dimexpr\textwidth-8\tabcolsep\relax}>{\raggedright\arraybackslash}p{0.25301\dimexpr\textwidth-8\tabcolsep\relax}>{\raggedright\arraybackslash}p{0.25301\dimexpr\textwidth-8\tabcolsep\relax}>{\raggedright\arraybackslash}X}
\toprule
\rowcolor{pale}
\textbf{指标} & \textbf{公开权重} & \textbf{累计 5 轮} & \textbf{变化（微调减基线）} \\
\midrule
{NDS} & {0.026860} & {0.118891} & {+0.092031} \\
{mAP} & {0.000061} & {0.062302} & {+0.062240} \\
\bottomrule
\end{tabularx}
\end{table}
\begin{table}[H]
\centering\small
\caption{微调前后检测误差}
\label{tab:8}
\begin{tabularx}{\textwidth}{>{\raggedright\arraybackslash}p{0.31325\dimexpr\textwidth-6\tabcolsep\relax}>{\raggedright\arraybackslash}p{0.34337\dimexpr\textwidth-6\tabcolsep\relax}>{\raggedright\arraybackslash}X}
\toprule
\rowcolor{pale}
\textbf{误差指标} & \textbf{公开权重} & \textbf{累计 5 轮} \\
\midrule
{mATE} & {1.15690} & {0.85550} \\
{mASE} & {0.90030} & {0.68340} \\
{mAOE} & {1.26350} & {0.95130} \\
{mAVE} & {1.68880} & {1.38360} \\
{mAAE} & {0.83150} & {0.63240} \\
\bottomrule
\end{tabularx}
\end{table}
{\small mAP 汇总不同中心距离阈值下的类别 AP；NDS 综合精度和位置、尺寸、朝向、速度、属性误差。mATE 单位米，mAOE 单位弧度，mAVE 单位米/\allowbreak{}秒。部分类别的误差项不适用或因匹配不足使用默认值，不宜解释为精确测得的单个目标误差。\par}

\begin{table}[H]
\centering\small
\caption{两个展示阈值下的预测框数量}
\label{tab:9}
\begin{tabularx}{\textwidth}{>{\raggedright\arraybackslash}p{0.31325\dimexpr\textwidth-6\tabcolsep\relax}>{\raggedright\arraybackslash}p{0.34337\dimexpr\textwidth-6\tabcolsep\relax}>{\raggedright\arraybackslash}X}
\toprule
\rowcolor{pale}
\textbf{固定验证帧} & \textbf{展示阈值 0.10} & \textbf{展示阈值 0.30} \\
\midrule
{帧 0} & {51} & {2} \\
{帧 1} & {45} & {4} \\
{帧 2} & {77} & {4} \\
\bottomrule
\end{tabularx}
\end{table}
阈值由 0.10 提高至 0.30，仅对缓存预测进行展示筛选，不重新运行模型，也不改变已计算的 NDS/\allowbreak{}mAP。BEV 图仅含原始点云和模型预测框，没有叠加真实框；框的数量或视觉密集程度不能替代检测评估。


本次微调的 NDS 变化为 +0.092031，mAP 变化为 +0.062240。结果仅对应 mini、单帧及 50×50 米教学范围，不能与完整 nuScenes 榜单直接比较；公开权重也不能视作从未见过相关数据的独立测试模型。



\FloatBarrier

\clearpage
\section{LiDAR BEV 结果展示}

\subsection{微调后的 LiDAR BEV：固定验证帧 0}

\begin{figure}[H]
\centering
\includegraphics[width=0.90\textwidth]{\detokenize{images/bev_01.png}}
\caption{累计五轮模型，展示阈值 0.10，共 51 个预测框。}
\label{fig:2}
\end{figure}
灰点为点云在 LiDAR x-y 平面的投影；彩色矩形表示预测三维框的底面和朝向，文字为类别及置信度。坐标单位为米，显示范围为 ±25 米。该视图未叠加真实框，需结合前述验证指标判断效果。


{\small 相同点云在阈值 0.30 下显示 2 个框。两套图片均保留，文件夹分别为 visualizations\_thr010 和 visualizations\_thr030。\par}

{\small 点云来源：n008-2018-08-01-15-16-36-0400\_\_LIDAR\_TOP\_\_1533151603547590.pcd.bin\par}


\FloatBarrier

\clearpage
\subsection{微调后的 LiDAR BEV：固定验证帧 1}

\begin{figure}[H]
\centering
\includegraphics[width=0.90\textwidth]{\detokenize{images/bev_02.png}}
\caption{累计五轮模型，展示阈值 0.10，共 45 个预测框。}
\label{fig:3}
\end{figure}
灰点为点云在 LiDAR x-y 平面的投影；彩色矩形表示预测三维框的底面和朝向，文字为类别及置信度。坐标单位为米，显示范围为 ±25 米。该视图未叠加真实框，需结合前述验证指标判断效果。


{\small 相同点云在阈值 0.30 下显示 4 个框。两套图片均保留，文件夹分别为 visualizations\_thr010 和 visualizations\_thr030。\par}

{\small 点云来源：n008-2018-08-01-15-16-36-0400\_\_LIDAR\_TOP\_\_1533151604048025.pcd.bin\par}


\FloatBarrier

\clearpage
\subsection{微调后的 LiDAR BEV：固定验证帧 2}

\begin{figure}[H]
\centering
\includegraphics[width=0.90\textwidth]{\detokenize{images/bev_03.png}}
\caption{累计五轮模型，展示阈值 0.10，共 77 个预测框。}
\label{fig:4}
\end{figure}
灰点为点云在 LiDAR x-y 平面的投影；彩色矩形表示预测三维框的底面和朝向，文字为类别及置信度。坐标单位为米，显示范围为 ±25 米。该视图未叠加真实框，需结合前述验证指标判断效果。


{\small 相同点云在阈值 0.30 下显示 4 个框。两套图片均保留，文件夹分别为 visualizations\_thr010 和 visualizations\_thr030。\par}

{\small 点云来源：n008-2018-08-01-15-16-36-0400\_\_LIDAR\_TOP\_\_1533151604547893.pcd.bin\par}

{\small 复现及证据：提交包中的 README.md 给出完整命令；manifest.json、comparison.json、final\_validation.json、训练审计、损失图、完整日志及 epoch\_5.pth 均随包提交。参考：课程指导书、随附 MMDetection3D 1.4.0 源码与 nuScenes SDK。\par}

\end{document}
